% === APPENDIX ================================================================
\section{Assumptions and Reproducibility Notes}
\label{subsec:assumptions}

The following are assumptions or design choices rather than measured facts.
They are listed explicitly so that a reader can judge how far each result
depends on them, and so that each can be discharged by a specific
measurement.

\subsection*{A1. Yaw inertia $I_{z}$}
$I_{z}=\SI{51.1}{\kilo\gram\meter\squared}$ is carried in the model file and
is \emph{not} measured. It enters \eqref{eq:omdot} directly and therefore
scales the yaw-rate residual the network is trained on, and it enters the
stiffness-adaptive sub-stepping of Section~\ref{subsec:gates} through
$l_{f}^{2}C_{f}+l_{r}^{2}C_{r}$ over $I_{z}$. \emph{How to discharge:} apply a
known yaw moment (a step steer at constant speed with the resulting axle
forces read from the force feedback topic) and fit $I_{z}$ from the
measured $\dot\omega$; or read the value the physics engine is using directly
from the actor's inertia tensor.

\subsection*{A2. Centre-of-gravity split}
$l_{f}=\SI{0.738142}{\meter}$, $l_{r}=\SI{0.795362}{\meter}$ are taken from
the controller configuration and confirmed by the vehicle owner. The measured
static load split (51.2/48.8 front/rear) implies a marginally different split
than the geometry does (51.9/48.1) --- a \SI{1.4}{\percent} disagreement. We
have not resolved it, and it is below the level at which any conclusion in
this paper would change.

\subsection*{A3. CG height and load transfer}
$h_{\mathrm{cg}}=\SI{0.3}{\meter}$ is carried in the model file and unused:
the single-track model of Section~\ref{subsec:model} neglects load transfer,
following the baseline. Normal loads are static. On a full-scale vehicle at
\SI{1}{\g} this is a stronger approximation than at 1:10, and it is a
candidate explanation for part of the residual the network absorbs.

\subsection*{A4. Steering-angle sign and units}
\texttt{feedback/\allowbreak steering\_angle} is published in degrees and is treated as
the front road-wheel angle, positive left, matching the ROS body frame. It was
verified to agree with the front-left wheel joint state to within
\SI{0.5}{\percent}. Left and right road wheels are not distinguished; the
single-track $\delta$ is taken as this single reported value.

\subsection*{A5. Steady-state force recovery}
Equation~\eqref{eq:ssforces} is not a force measurement --- it is fully
determined by $v_{x}\omega$ --- so it misattributes transients. It is applied
only to the \emph{synthetic}, quasi-steady rollout, where the steady-state
assumption is constructed rather than hoped for. Its adequacy was nonetheless
checked on measured data against the simulator's own axle force
(correlation \num{0.993}, slope \num{0.94}); we treat that as validating the
relation, not as validating a steady-state assumption on recorded laps.

\subsection*{A6. The residual network is not the plant}
The corrected model $\hat\vecx_{k+1}+e_{k}$ is queried by the rollout at
inputs the network has only partially seen. The extrapolation bound of
Section~\ref{subsec:iterative} limits but does not eliminate this: the sweep
still reaches $1.5\times$ the observed 99th percentile of $|\delta|$. The
choice of $1.5$ is empirical, selected because it was sufficient to reach the
tire's peak at the corrected rollout speed while removing the observed drift;
it is not derived.

\subsection*{A7. Simulator determinism and timing}
CARLA runs in synchronous mode with a \SI{0.033}{\second} physics step, but
sim time advances faster than wall time (observed $\approx1.9\times$ on our
hardware). All identification data is stamped from the simulation clock.

\subsection*{A8. Software configuration of record}
Identification and control workspace: ROS~2 Humble, Python 3.10 (required ---
ROS~2 Humble ships CPython~3.10 ABI extensions only), PyTorch with CUDA on an
RTX~2080-class GPU. The bridge is built and sourced as a separate overlay on
top of that workspace, so that its message packages
(\texttt{sim\_manager\_\allowbreak msgs}) resolve for the ground-truth forces;
CARLA server 0.9.16 / Unreal Engine 4.26; \textbf{acados 0.5.5} with
partial-condensing HPIPM. The identifier's acquisition and state-machine timer
runs at \SI{50}{\hertz} (\SI{20}{\hertz} is the controller's period) with a
\SI{30}{\second} buffer, six co-identification iterations, 200 training
epochs, full-batch Adam at $5\times10^{-4}$. The one-step model is discretised
at $T_{s}=\SI{0.035}{\second}$. This is
an empirical value---it matches neither the server's \SI{0.033}{\second} step
nor the \SI{50}{\hertz} acquisition timer---and is carried because the
identification is insensitive to it over the range tested
(Section~\ref{subsec:data}).

\emph{Random seeds.} \texttt{nn\_train} seeds PyTorch and NumPy at \num{42} for
the single-member architectures used here (\num{42}$+i$ per ensemble member).
\texttt{pacejka\_solver.seed} is \texttt{null} in both shipped parameter sets,
so the multi-start jitter and---for \texttt{Ours}---the differential-evolution
population draw from an unseeded generator: the reported identifications are
reproducible in distribution but not bit for bit. The sweep of
Table~\ref{tab:sweep} is seeded explicitly and is bit-reproducible.

\section{Simulator--ROS~2 Interface}
\label{sec:iface}

\subsection*{Command path}
The bridge runs in \texttt{ackermann\_drive} mode, subscribing to a single
\texttt{ack\-er\-mann\_msgs/AckermannDriveStamped} on \texttt{/drive} and
forwarding speed and steering angle to CARLA's native Ackermann controller,
whose server-side cascaded PID loops are configured once at startup. No
client-side steering PID is present---the commanded road-wheel angle passes
through in radians---so the achieved/\allowbreak commanded ratio of
Table~\ref{tab:vehicle} is a property of the vehicle and its steering curve
rather than of a tuning choice in the bridge.

\subsection*{Feedback path}
Table~\ref{tab:topics} lists the interface consumed by the identifier and its
validation. One property of it is specific to this architecture.

\begin{table}[t]
\caption{ROS~2 interface used by the identification and validation pipeline.
Namespace is \texttt{/sim} unless marked absolute.}
\label{tab:topics}
\centering
\scriptsize
\setlength{\tabcolsep}{3pt}
\begin{tabular}{@{}p{2.4cm}p{2.75cm}p{2.5cm}@{}}
\toprule
Topic & Type & Role\\
\midrule
\texttt{/odom} & \texttt{nav\_msgs/Odometry} & State $v_{x},v_{y},\omega$; twist is body-frame\\
\texttt{feedback/\allowbreak steering\_angle} & \texttt{std\_msgs/Float32} (deg) & Achieved road-wheel angle $\delta$\\
\texttt{feedback/imu} & \texttt{sensor\_msgs/Imu} & $a_{x},a_{y}$ for the friction warm-start\\
\texttt{feedback/\allowbreak tire\_forces} & \texttt{sim\_manager\_\allowbreak msgs/\allowbreak TireForces} & Per-wheel $\alpha$, $F_{z}$, $F_{y}$, $\mu$, $\omega_{w}$, $\kappa$: \textbf{ground truth}\\
\texttt{feedback/\allowbreak vehicle\_physics} & \texttt{std\_msgs/String} (JSON) & Live wheel parameters of \eqref{eq:physx}\\
\texttt{feedback/speed} & \texttt{std\_msgs/Float32} & Forward speed echo\\
\texttt{feedback/motors} & \texttt{std\_msgs/String} (JSON) & Per-wheel speed/torque/brake\\
\texttt{/drive} & \texttt{ackermann\_msgs/\allowbreak Ackermann\allowbreak Drive\allowbreak Stamped} & Control command in\\
\texttt{control/\allowbreak tire\_friction} & \texttt{std\_msgs/Float32} & Runtime grip override\\
\texttt{/clock} \emph{(abs.)} & \texttt{rosgraph\_msgs/\allowbreak Clock} & Simulation time base\\
\bottomrule
\end{tabular}
\end{table}

\paragraph{Per-wheel ground truth}
\texttt{feedback/\allowbreak tire\_forces} carries CARLA's own
\texttt{Wheel\-TelemetryData}: per-wheel lateral slip angle, normal load,
lateral force, effective friction coefficient, wheel speed, longitudinal force
and wheel torque, in wheel order [FL, FR, RL, RR]. No analytic tire model
exists inside the bridge, so nothing on this topic can disagree with the
physics the vehicle is driven by. This topic is what allows an on-track
identifier to be scored against the forces the plant generated. Four properties
of it were established by measurement and bound what may be read off it.
(i)~Only \texttt{lateral\_force} is force ground truth; it is the channel every
validation in this paper uses. (ii)~\texttt{longitudinal\_force} is a
friction-capacity report rather than a measurement---it tracks
$\mu F_{z}$ and correlates with $ma_{x}$ at \num{0.04}---and
\texttt{wheel\_torque} is an exact identity of it, so neither carries
independent information. (iii)~Below \SI{0.5}{\meter\per\second} the publisher
zeroes slip angle and both forces while the normal load stays live, so a
load-only gate admits standstill frames as perfect $(0,0)$ samples; the
identifier gates on the all-zero pattern as well. (iv)~The slip angle needs no
sign change across the boundary, which was verified rather than assumed.

Figure~\ref{fig:frames} states the conventions the rest of the paper assumes.
The simulator and the identification stack do not share a handedness: CARLA
inherits Unreal's left-handed world (\mbox{$X$ east}, $Y$ \emph{south}, $Z$ up,
yaw positive clockwise), while every ROS~2 topic follows
REP-103~\cite{rep103}---right-handed ENU for the world, forward-left-up for the
body---and the identification equations of Section~\ref{sec:method} are written
in the body frame with $\alpha$ and $F_{y}$ positive to the left, the SAE-style
axis convention~\cite{saej670,pacejka2012book}. The bridge is where the two
meet, and it is a per-channel conversion rather than a single transform:
position and yaw flip sign, twist is additionally rotated into the body frame
by $q^{-1}$, the IMU's $a_{y}$ flips, the steering \emph{command} flips on the
way out, and the tire telemetry's slip angle does \emph{not}, because the frame
flip and CARLA's own sign convention for \texttt{lat\_slip} cancel. That last
one is the trap: a sign error there is invisible in the states and inverts the
identified curve, so it was checked against a slip angle rebuilt from the
chassis state in a steady \SI{4}{\meter\per\second} turn at
\SI{0.20}{\radian}---equal magnitude to within \SI{2}{\percent} front and
\SI{15}{\percent} rear.

\begin{figure*}[t]
\centering
\resizebox{\textwidth}{!}{% =============================================================================
% Fig. -- Frame and sign conventions across the CARLA/ROS 2 boundary.
%
% Vector TikZ, drawn at the paper's own font sizes. Included by
% sections/platform.tex inside a figure* (full text width).
%
% Three panels, left to right:
%   (a) the two world frames and the map between them (CARLA/Unreal is
%       left-handed with Y south and yaw clockwise; ROS is ENU/REP-103).
%   (b) the vehicle body frame in which every identification quantity is
%       defined: FLU, positive steer to the left, positive lateral force to
%       the left, slip angles as in (2).
%   (c) the per-channel conversion the bridge applies, which is where a sign
%       error would enter the identification silently.
%
% Layout note: every panel's bounding box is fitted to two explicit corner
% coordinates rather than to its content, so the three panels keep the same
% height and nothing collides once the whole picture is resized to \textwidth.
% =============================================================================

\begin{tikzpicture}[
  font=\scriptsize,
  >={Stealth[length=1.7mm,width=1.5mm]},
  ax/.style      = {->, line width=0.5pt, black!75},
  car/.style     = {line width=0.7pt, black!75},
  wheel/.style   = {line width=1.5pt, black!80, line cap=round},
  frc/.style     = {->, line width=0.7pt, green!45!black},
  vel/.style     = {->, line width=0.6pt, blue!55!black},
  rot/.style     = {->, line width=0.5pt, black!60},
  lbl/.style     = {font=\tiny, inner sep=1.0pt},
  note/.style    = {font=\tiny, align=left, inner sep=1.6pt},
  hdr/.style     = {font=\scriptsize\bfseries, align=center},
  conv/.style    = {->, line width=0.6pt, black!65, dashed}
]

% ================== (a) vehicle body frame and sign rules ==================
% Single-track sketch: x forward (right), y left (up), omega CCW positive.
\begin{scope}[shift={(93mm,2mm)}]
  \coordinate (cg)  at (0,0);
  \coordinate (fr)  at (15mm,0);
  \coordinate (rr)  at (-13mm,0);

  \draw[car] (rr) -- (fr);
  \draw[wheel] ($(rr)+(-3mm,0)$) -- ($(rr)+(3mm,0)$);
  \draw[wheel, rotate around={17:(fr)}] ($(fr)+(-3mm,0)$) -- ($(fr)+(3mm,0)$);
  \fill (cg) circle (0.6mm);

  % body axes at the centre of gravity
  \draw[ax] (cg) -- ++(8mm,0) node[lbl,below] {$x$ fwd};
  \draw[ax] (cg) -- ++(0,7mm) node[lbl,left]  {$y$ left};
  \draw[rot] (-2.8mm,-2.8mm) arc (225:135:4.0mm);
  \node[lbl, anchor=north east] at (-3.2mm,-2.4mm) {$\omega>0$};

  % steer angle at the front wheel
  \draw[black!45, line width=0.35pt] (fr) -- ++(6mm,0);
  \draw[rot] ($(fr)+(5mm,0)$) arc (0:17:5mm);
  \node[lbl, anchor=west] at ($(fr)+(5.4mm,2.2mm)$) {$\delta>0$};

  % velocity and sideslip at the CG
  \draw[vel] (cg) -- ++(10mm,4.5mm) node[lbl,above] {$v$};
  \draw[rot] (5.5mm,0) arc (0:24.2:5.5mm);
  \node[lbl, anchor=west] at (6.0mm,1.2mm) {$\beta$};

  % axle lateral forces, positive to the left
  \draw[frc] (fr) -- ++(0,7mm) node[lbl,above] {$F_{y,f}>0$};
  \draw[frc] (rr) -- ++(0,7mm) node[lbl,above] {$F_{y,r}>0$};
\end{scope}

\node[note, anchor=north west, text width=62mm] at (63mm,-11mm)
  {$\alpha_{f}=\delta-\arctan\dfrac{v_{y}+l_{f}\omega}{v_{x}}$,\qquad
   $\alpha_{r}=-\arctan\dfrac{v_{y}-l_{r}\omega}{v_{x}}$\\[2pt]
   Positive $\alpha$ produces positive (leftward) $F_{y}$. Wheel order in
   every per-wheel array: [FL, FR, RL, RR].};

\coordinate (pb1) at (61mm,-27mm);
\coordinate (pb2) at (127mm,15mm);
\begin{scope}[on background layer]
  \node[draw, dashed, rounded corners=2pt, inner sep=0pt, fill=black!2,
        fit=(pb1)(pb2)] (panelb) {};
\end{scope}
\node[hdr, above=1.0mm of panelb.north] {(a) Vehicle body frame (FLU)};

% ====================== (b) per-channel conversions =========================
\node[note, anchor=north west] at (133mm,13mm)
  {\renewcommand{\arraystretch}{1.2}%
   \begin{tabular}{@{}p{17mm}p{40mm}@{}}
   \textbf{Channel} & \textbf{Bridge conversion}\\
   \hline\\[-6pt]
   \texttt{/odom} pose & $y\mapsto-y$; $\psi,\theta\mapsto-\psi,-\theta$\\
   \texttt{/odom} twist & world $\to$ body by $q^{-1}$; $v_{y}$,
   $\omega_{y}$, $\omega_{z}$ negated\\
   \texttt{feedback/imu} & $a_{y}\mapsto-a_{y}$\\
   \texttt{tire\_forces} $\alpha$ & $\deg\to\mathrm{rad}$ only: the frame flip
   and CARLA's own sign convention cancel\\
   \texttt{tire\_forces} $F_{y}$ & $F_{y}\mapsto-F_{y}$, into $+y$ left\\
   \texttt{/drive} $\delta$ & negated on the way out to CARLA\\
   \end{tabular}};

\coordinate (pc1) at (131mm,-27mm);
\coordinate (pc2) at (196mm,15mm);
\begin{scope}[on background layer]
  \node[draw, dashed, rounded corners=2pt, inner sep=0pt, fill=black!2,
        fit=(pc1)(pc2)] (panelc) {};
\end{scope}
\node[hdr, above=1.0mm of panelc.north] {(b) Conversions applied by \bridge{}};

\end{tikzpicture}
}
\caption{Frame and sign conventions across the simulator--ROS~2 boundary.
(a) All identification quantities are defined in the body frame:
$x$ forward, $y$ left, $\omega$ and $\delta$ positive counter-clockwise, and
slip angle and lateral force positive to the left, so that a positive $\alpha$
produces a positive $F_{y}$~\cite{saej670}. (b) The conversion applied per
channel. The slip angle of the tire telemetry is the one exception that needs
no sign change, because the frame flip and CARLA's own convention cancel.}
\label{fig:frames}
\end{figure*}

\subsection*{Extensions made for identification}
Six changes support this study and are listed for reproducibility: the
existing but unsubscribed \texttt{feedback/\allowbreak steering\_angle} is now
consumed by the identifier (with configurable units, scale and staleness
timeout, falling back to the \texttt{/drive} setpoint); runtime overrides
\texttt{control/\allowbreak tire\_friction} and
\texttt{control/\allowbreak drag\_coefficient} write straight to the live
\texttt{VehiclePhysicsControl}, so surface friction can change within a run
without respawning; the torque curve, zero-throttle damping and wheel
\texttt{tire\_friction} were tuned against a low-speed hold test (a wheel
friction above $\approx$2.0 on a steered wheel scrubs hard enough at creep
speed to stall the drivetrain intermittently), and these values define the
plant being identified; telemetry runs on a dedicated \SI{10}{\hertz}
thread off the control loop, since each tick costs several blocking CARLA RPCs;
\texttt{feedback/\allowbreak tire\_forces} was extended with the effective
friction coefficient, the wheel speed and a kinematically recomputed slip ratio
$\kappa=(\omega_{w}r_{w}-v)/v$, the last because CARLA's own
\texttt{long\_slip} field is inconsistent with the wheel speed and radius the
same message reports (\num{+0.177}, i.e.\ hard braking, on a wheel turning
\SI{8.99}{\radian\per\second} at a steady \SI{2.33}{\meter\per\second} cruise,
whose kinematic ratio is \num{-0.034}); and the live wheel parameters of
\eqref{eq:physx} are published on
\texttt{feedback/\allowbreak vehicle\_physics}, so the reference curve of
Table~\ref{tab:ref} cannot go stale when the vehicle configuration changes.

\section{Reference Implementation}
\label{sec:refimpl}

The identification package, the ROS~2 bridge and the controller stack used in
this paper are available as \texttt{On-Track-SysID}, extended
from~\cite{ontrack2025} (upstream:
\url{https://github.com/ForzaETH/On-Track-SysID}), \bridge{}
(\url{https://github.com/ebrahimabdelghfar/Carla_ASU_Bridge})%
~\cite{carlaasubridge}, and the accompanying MPC and benchmarking packages,
which are released together with the scenario definitions, the exported result
CSVs and the figure and sweep generators at
\url{https://github.com/ebrahimabdelghfar/Ebrahim_Master_Thesis_repo/tree/feat-follow-premade-racing-track}. The benchmarking packages compute one-step and multi-step
prediction metrics online (RMSE, MAE, NRMSE, MaxAE, bias, standard deviation
and $R^{2}$ by Welford's algorithm), compare estimated against ground-truth
per-wheel lateral forces, and record controller-switching behaviour, and were
built for this study.
